\section{A long prime polynomial}\label{lp:section}

The correlation argument will require cancellation in a prime polynomial at
frequencies as large as $x^2$.  We prove the required estimate here, including
the dependence on a growing degree in the elementary mean-value argument.
Throughout this section, $L=\log x$ and $T=\log L$.

\begin{lemma}[Long prime polynomial]\label{lp:prime-polynomial}
Fix $0<\tau<\eta<1$, $C>0$, and $A>0$.  There is a constant
$B_0=B_0(\tau,\eta,C,A)$ such that, uniformly for
\[
 \frac{x^\tau}{2}\le N\le x^\eta,\qquad
 q\le L^C,\qquad L^{B_0}\le |t|\le x^2,
\]
every Dirichlet character $\chi$ modulo $q$ and every interval
$I\subset[N,2N]$ satisfy
\begin{equation}\label{lp:prime-estimate}
 \left|\sum_{p\in I}\chi(p)p^{-1+it}\right|
 \ll_{\tau,\eta,C,A} L^{-A}.
\end{equation}
\end{lemma}

The proof proceeds from a quantitative power-sum estimate to cancellation
for a logarithmic phase, then to a high-height zero-free strip, and finally
to primes by Mellin inversion.
We first establish a power-sum estimate with sufficient uniformity in its
degree.  For integers $k\ge2$, $s\ge1$, and $M\ge1$, let $J_{s,k}(M)$ count
the solutions of
\[
 \sum_{i=1}^s u_i^j=\sum_{i=1}^s v_i^j
 \quad(1\le j\le k),\qquad 1\le u_i,v_i\le M.
\]
Writing $K=k(k+1)/2$ and
\[
 f(\boldsymbol\alpha)=\sum_{n=1}^M
       \e\!\left(\sum_{j=1}^k\alpha_j n^j\right),
\]
orthogonality gives
$J_{s,k}(M)=\int_{[0,1]^k}|f(\boldsymbol\alpha)|^{2s}
\,d\boldsymbol\alpha$.

The argument uses the classical Vinogradov mean-value method in
Linnik's $p$-adic form; see Wooley~\cite[Section~2, pp.~1583--1585]{Wooley2012}
for an exposition.  We derive the required degree-uniform quantitative
estimate below, without invoking the modern efficient-congruencing
theorem of that paper.

\begin{lemma}[A quantitative power-sum bound]\label{lp:mean-value}
There are absolute constants $C_1,C_2>0$ such that, for every integer
$k\ge2$, some integer $s$ with $k\le s\le C_1k^4$ satisfies
\begin{equation}\label{lp:mean-value-bound}
 J_{s,k}(M)\le \exp(C_1k^{C_2})M^{2s-K+1/100}
 \qquad(M\ge1).
\end{equation}
\end{lemma}

\begin{proof}
We iterate estimates
\begin{equation}\label{lp:inductive-bound}
 J_{s,k}(M)\le C_s M^{E_s},\qquad
 E_s=2s-K+\epsilon_s,
\end{equation}
starting with $s=k$, $C_k=1$, and $\epsilon_k=K$.  The iteration sends
$s$ to $s+k$ and $\epsilon_s$ to $(1-1/k)\epsilon_s$.

Choose a sufficiently large absolute constant $A_1$.  If
$M^{1/k}<A_1k^6$, then $\log M\ll k\log(2k)$, and the trivial bound
$J_{u,k}(M)\le M^{2u}$ costs at most
\begin{equation}\label{lp:small-M-cost}
 M^K\le\exp\bigl(O(k^3\log(2k))\bigr)
\end{equation}
relative to every proposed estimate with exponent $2u-K+\epsilon$,
$\epsilon\ge0$.  Thus it suffices to treat $M^{1/k}\ge A_1k^6$.
The prime number theorem in Proposition~\ref{pre:primes}, after enlarging
$A_1$, supplies a fixed list of
$2k^3+5$ primes $r$ in $[M^{1/k},2M^{1/k}]$, all greater than $k$.

At moment $s+k$, call a tuple degenerate if it has fewer than $k$
distinct coordinates.  There are at most
$k^{s+k}M^{k-1}$ such tuples.  If $G$ is their exponential sum, its
contribution on one side of the equations is at most
\[
 \int_{[0,1]^k}|G|\,|f|^{s+k}
 \le k^{s+k}M^{k-1}J_{s+k,k}(M)^{1/2}.
\]
The contribution with a degenerate tuple on either side is therefore at
most twice this quantity.

For a nondegenerate solution, select $k$ distinct coordinates on each side
and permute them to the first $k$ positions.  This costs at most
$(s+k)^{2k}$.  The product of the two Vandermonde products is a nonzero
integer of absolute value at most $M^{k(k-1)}$.  Fewer than $k^2(k-1)+1$
primes of size at least $M^{1/k}$ can divide it.  Hence some prime $r$ in
our list makes these first $k$ coordinates distinct modulo $r$ on each
side.

For such a prime put
\[
 F_r(\boldsymbol\alpha)=
 \sum_{\substack{1\le z_1,\ldots,z_k\le M\\
                  z_i\not\equiv z_j\pmod r\ (i\ne j)}}
 \e\!\left(\sum_{j=1}^k\alpha_j\sum_{i=1}^kz_i^j\right).
\]
The relevant number of solutions is bounded by
\[
 (s+k)^{2k}\sum_r\int_{[0,1]^k}|F_r|^2|f|^{2s}.
\]
The integrands are nonnegative.  Write $f=\sum_{a\bmod r}f_a$, where
$f_a$ is restricted to $n\equiv a\pmod r$.  H\"older's inequality gives
\[
 |f|^{2s}\le r^{2s-1}\sum_{a\bmod r}|f_a|^{2s}.
\]

Fix $a$.  In the system counted by
$\int|F_r|^2|f_a|^{2s}$, translate all variables by $-a$.
Translation preserves the equations for the first $k$ powers by the
binomial formula.  The remaining $s$ variables on either side are
multiples of $r$, so the first lists $\boldsymbol z,\boldsymbol z'$ obey
\begin{equation}\label{lp:power-congruences}
 \sum_{i=1}^kz_i^j\equiv\sum_{i=1}^k(z_i')^j\pmod{r^j}
 \qquad(1\le j\le k).
\end{equation}
There are at most $M^k$ choices for the first list.  For each such list,
there are at most $k!r^{K-k}$ choices for the second.  To see this,
lift the prescribed $j$th sum modulo $r^j$ to a residue modulo $r^k$.
The number of choices for all lifts is
$\prod_{j=1}^k r^{k-j}=r^{K-k}$.
For each full vector of sums modulo $r^k$, Newton's identities determine
the multiset of roots modulo $r$, since $r>k$.  There are at most $k!$
orderings.  The Jacobian of the power sums has determinant
\[
 k!\prod_{i<j}(z_j'-z_i'),
\]
up to sign, and is invertible modulo $r$.  Each ordering therefore lifts
uniquely from modulus $r$ to modulus $r^k$: at each stage the next digits
are the unique solution of the corresponding linear system modulo $r$.
Finally, an interval of $M$ integers contains at most one representative
of any residue modulo $r^k$, because $M\le r^k$.

After both first lists are fixed, divide the remaining variables by $r$.
They range over a common interval of at most $\lceil M/r\rceil$ integers
and have prescribed differences of power sums.  Translation to an
initial interval changes only these prescribed differences.  The count
is a Fourier coefficient of the nonnegative function $|f|^{2s}$ at that
shorter length, and consequently is at most
$J_{s,k}(\lceil M/r\rceil)$.  Summing over $a$ accounts for the final
factor $r$, and gives
\begin{equation}\label{lp:recurrence-count}
 \begin{split}
 J_{s+k,k}(M)
 &\le 2k^{s+k}M^{k-1}J_{s+k,k}(M)^{1/2}\\
 &\quad +(s+k)^{2k}(2k^3+5)k!
       \max_r r^{2s+K-k}M^kJ_{s,k}(\lceil M/r\rceil).
 \end{split}
\end{equation}

Insert \eqref{lp:inductive-bound}.  Since $E_s\ge0$, rounding costs
at most $2^{E_s}$, and the exponent of $r$ becomes
\[
 2s+K-k-E_s=k^2-\epsilon_s\ge0.
\]
Replacing $r$ by at most $2M^{1/k}$ therefore gives the exponent
\[
 k+E_s+\frac{k^2-\epsilon_s}{k}
 =2(s+k)-K+(1-1/k)\epsilon_s.
\]
The inequality $J\le a\sqrt J+b$ implies $J\le2a^2+2b$.
The exponent $2k-2$ from $a^2$ is admissible: the target exponent
starts at $2k$, and at each step its increase is
$2k-\epsilon_s/k\ge2k-K/k>0$.
Thus the asserted iteration holds.  Its constants can be chosen with
\[
 \log C_{s+k}\le \log(1+C_s)
 +O\!\left((s+k)\log(2k)+k\log(s+k)+k^3\log(2k)\right),
\]
including \eqref{lp:small-M-cost}.

After $O(k\log(2k))$ steps,
$\epsilon_s=K(1-1/k)^{(s-k)/k}\le1/100$.
Then $s=O(k^2\log(2k))\le C_1k^4$, and summing the displayed costs
gives $\log C_s\le C_1k^{C_2}$ for absolute constants.
Increasing the exponent from $\epsilon_s$ to $1/100$ proves the Lemma.
\end{proof}

\begin{lemma}[A logarithmic phase on progressions]\label{lp:log-phase}
Fix $C>0$.  There is an absolute constant $C_3>0$ such that, for
sufficiently large $x$ in terms of $C$, the following holds uniformly:
\[
 \exp(L/T^2)\le N'\le2x^5,\qquad q\le L^C,\qquad
 \exp(L/(2T^2))\le |v|\le4x^3.
\]
For every residue $a\pmod q$ and every interval $I\subset[N',2N']$,
\begin{equation}\label{lp:progression-estimate}
 \left|\sum_{\substack{n\in I\\n\equiv a\pmod q}}n^{iv}\right|
 \ll \frac{N'}q\exp(-L/T^{C_3}).
\end{equation}
\end{lemma}

\begin{proof}
Put $H=N'/q$ and $y=\log|v|/\log H$.  Then
\begin{equation}\label{lp:phase-scales}
 \log H\ge L/T^2-CT\gg L/T^2,\qquad y\ll T^2.
\end{equation}
Writing $n=q(b+a/q)$, with $0\le a<q$, reduces the sum, up to a
factor of modulus one, to $\sum_{b\in\mathcal I}(b+a/q)^{iv}$,
where $\mathcal I$ is an interval of integers and
$H\le b+a/q\le2H$.

If $y<4/5$, the derivative of
$v\log(b+a/q)/(2\pi)$ is monotone, has magnitude comparable to
$|v|/H$, and has magnitude less than $1/2$.  The monotone first
derivative estimate in Proposition~\ref{pre:derivatives} therefore bounds
the sum by $O(H/|v|)$.
The lower bound on $|v|$ makes this smaller than
\eqref{lp:progression-estimate}.

Suppose now that $y\ge4/5$.  Set
\[
 M=\lfloor H^{3/4}\rfloor,\qquad k=\lceil4y+8\rceil.
\]
Averaging the sum over forward shifts $1\le h\le M$ changes it by
$O(M)$: a shift changes an interval at only $O(h)$ endpoints.
For $b\in\mathcal I$, Taylor expansion gives
\[
 \frac{v}{2\pi}\log(b+h+a/q)
 =\frac{v}{2\pi}\log(b+a/q)
       +\sum_{j=1}^k\alpha_j(b)h^j+O(H^{-2}),
 \qquad
 \alpha_j(b)=\frac{(-1)^{j-1}v}{2\pi j(b+a/q)^j}.
\]
Indeed the remainder is
$O(|v|(M/H)^{k+1})=O(H^{y-(k+1)/4})=O(H^{-9/4})$,
uniformly even as $k$ grows.  Thus, with
\[
 S(\boldsymbol\alpha)=\sum_{h=1}^M
                 \e\!\left(\sum_{j=1}^k\alpha_jh^j\right),
\]
the original sum is bounded by
\begin{equation}\label{lp:shifted-phase}
 \frac1M\sum_{b\in\mathcal I}|S(\boldsymbol\alpha(b))|
       +O(M+H^{-1}).
\end{equation}

Partition the coefficient torus $[0,1)^k$ into boxes with side length
$M^{-j}$ in coordinate $j$.  Each box contains at most
\begin{equation}\label{lp:occupancy}
 \exp(O(k))H^{4/5}
\end{equation}
of the vectors $\boldsymbol\alpha(b)\pmod1$.
For this it suffices to use the coordinate
$j=\lceil y+3/20\rceil$, which lies between $1$ and $k$.
Before reduction modulo one, this coordinate has magnitude
$O(H^{-3/20})$, is monotone, and has derivative of magnitude at least
\[
 \frac{|v|}{2\pi(2H)^{j+1}}
       \ge\exp(-O(k))H^{y-j-1}.
\]
A coordinate interval of length $M^{-j}$ on the torus pulls back to
at most two intervals in $b$.  Their total length is at most
$\exp(O(k))H^{1+j/4-y}$.  Here the floor in $M$ costs
$\exp(O(k))$, while
\[
 y-\frac j4\ge\frac{3y}4-\frac{23}{80}
       \ge\frac5{16}>\frac15.
\]
Counting integer points proves \eqref{lp:occupancy}.

Let $s$ be supplied by Lemma~\ref{lp:mean-value}.  We also need the
following bound for suprema over boxes $Q$:
\begin{equation}\label{lp:box-supremum}
 \sum_Q\sup_{\boldsymbol\alpha\in Q}|S(\boldsymbol\alpha)|^{2s}
 \le\exp(O(sk+k))M^KJ_{s,k}(M).
\end{equation}
To prove it, rescale each box to $[0,1]^k$.  Iterating the
one-dimensional fundamental theorem of calculus gives, for any smooth
function $F$ on that cube,
\[
 \sup|F|\le\sum_{\boldsymbol\epsilon\in\{0,1\}^k}
       \int_{[0,1]^k}|\partial^{\boldsymbol\epsilon}F|.
\]
H\"older's inequality bounds the $2s$th power of the right-hand side
by $2^{k(2s-1)}$ times the sum of the corresponding $2s$th moments.
For the rescaled $S$, every mixed derivative has coefficients bounded
in modulus by $(2\pi)^k$: each differentiation in coordinate $j$
introduces the factor $2\pi i h^j/M^j$.
On summing over the boxes, change of variables contributes $M^K$.
Orthogonality then bounds the full-torus moment of every such derivative
by $(2\pi)^{2sk}J_{s,k}(M)$.  This proves
\eqref{lp:box-supremum}, with constants controlled at the growing degree.

Apply H\"older's inequality to the sum in
\eqref{lp:shifted-phase}, and use \eqref{lp:occupancy},
\eqref{lp:box-supremum}, and \eqref{lp:mean-value-bound}.  Since
$|\mathcal I|\ll H$, the result is
\begin{equation}\label{lp:phase-final-bound}
 \frac1M\sum_{b\in\mathcal I}|S(\boldsymbol\alpha(b))|
 \ll H\left(\exp(O(k^{C_4}))H^{-1/5}M^{1/100}\right)^{1/(2s)}
\end{equation}
for an absolute constant $C_4$.
We have $k\ll T^2$ and $s\ll T^8$, whereas
\[
 -\frac15\log H+\frac1{100}\log M
       \le-\frac{77}{400}\log H.
\]
Every fixed power of $T$ is $o(L/T^2)$.  Thus the right-hand side
of \eqref{lp:phase-final-bound} is at most
$H\exp(-c_1L/T^{10})$ for some absolute $c_1>0$, once $x$ is
sufficiently large.  The errors in \eqref{lp:shifted-phase} are smaller.
Increasing a fixed exponent $C_3>10$ absorbs $c_1$ and proves the Lemma.
\end{proof}

We write $D(s,\chi)$ for the Dirichlet $L$-function, to distinguish it
from $L=\log x$.  Characters need not be primitive.  Periodicity gives
$\sum_{n\le u}\chi(n)=c_\chi u+O(q)$, where
$c_\chi=q^{-1}\sum_{a\bmod q}\chi(a)$.  Partial summation therefore
continues $D(s,\chi)$ meromorphically to $\operatorname{Re}s>0$, with
only the possible simple pole at $s=1$, and gives, for
$\sigma=\operatorname{Re}s$ in a fixed compact subinterval of $(0,\infty)$,
\begin{equation}\label{lp:L-tail}
 D(s,\chi)=\sum_{n\le Y}\frac{\chi(n)}{n^s}
       +\frac{c_\chi Y^{1-s}}{s-1}
       +O\bigl(q(1+|s|)Y^{-\sigma}\bigr).
\end{equation}

\begin{lemma}[A bound near the line $\operatorname{Re}s=1$]
\label{lp:L-upper}
Fix $C>0$ and put $r_*=T^4/L$.  Uniformly for $q\le L^C$,
\begin{equation}\label{lp:L-upper-bound}
 \log|D(\sigma+iv,\chi)|\ll_C T^2
 \quad\left(|\sigma-1|\le10r_*,\quad
                   \frac12\le|v|\le3x^3\right).
\end{equation}
\end{lemma}

\begin{proof}
First suppose $|v|\le\exp(L/(2T^2))$, and use
$Y=\exp(2L/T^2)$ in \eqref{lp:L-tail}.  Absolute summation gives
\[
 \sum_{n\le Y}n^{-\sigma}
 \ll(1+\log Y)\max(1,Y^{1-\sigma})\le\exp(O(T^2)).
\]
Since $|s-1|\ge1/2$, the pole term has the same bound.  The error is
at most
\[
 \exp\left(-\frac{3L}{2T^2}+O(T^2+CT)\right),
\]
and hence is negligible.

For $\exp(L/(2T^2))<|v|\le3x^3$, take $Y=x^5$.
The terms with $n\le\exp(L/T^2)$ again contribute
$\exp(O(T^2))$ in absolute value.  On a dyadic interval above this
threshold, sum \eqref{lp:progression-estimate}, with phase $-v$, over
the residue classes modulo $q$, including their character coefficients.
The factors $q$ and $1/q$ cancel.  Partial summation with $n^{-\sigma}$
bounds that dyad by
\[
 \exp(-L/T^{C_3})\exp(O(T^4)).
\]
There are $O(L)$ dyads, and the same bound applies to a final partial
dyad.  Their total is negligible, since $L/T^{C_3}$ exceeds every
fixed power of $T$.  Finally, the pole term and remainder in
\eqref{lp:L-tail} are bounded respectively by
\[
 \exp\left(-\frac{L}{2T^2}+O(T^4)\right),\qquad
 \exp\bigl(-2L+O(T^4+CT)\bigr).
\]
This proves \eqref{lp:L-upper-bound}.
\end{proof}

\begin{lemma}[Local logarithmic derivative]\label{lp:local-log-derivative}
Fix $C>0$, let $q\le L^C$, and put
$s_0=1+r_*+iv$, where $1\le|v|\le\tfrac52x^3$.
There is a radius $R$ with $4r_*\le R\le5r_*$, with no zero on
its boundary, for which
\begin{equation}\label{lp:local-log-derivative-formula}
 \frac{D'}{D}(s,\chi)
 =\sum_{|\rho-s_0|<R}\frac1{s-\rho}
       +O_C(T^2/r_*)\qquad(|s-s_0|\le2r_*),
\end{equation}
away from zeros.  The sum counts zeros with multiplicity and has
$O_C(T^2)$ terms.
\end{lemma}

\begin{proof}
For large $x$, the disk $|s-s_0|\le8r_*$ avoids the possible pole at
$s=1$ and lies in the region of Lemma~\ref{lp:L-upper}.
At its center the Euler product gives
\[
 |D(s_0,\chi)|\ge\prod_p(1+p^{-1-r_*})^{-1}
       =\frac{\zeta(2+2r_*)}{\zeta(1+r_*)}\gg r_*.
\]
Thus $\log|D(s_0,\chi)|\ge-O(T)$.  Jensen's formula, using the
radius $8r_*$ and the upper bound $O_C(T^2)$, shows that the disk
of radius $5r_*$ contains $O_C(T^2)$ zeros.  Choose $R\in[4r_*,5r_*]$
so that none lies on its boundary.

Use centered coordinates $z=s-s_0$, and write $a=\rho-s_0$ for
each zero in $|z|<R$.  Divide $D(s_0+z,\chi)$ by the disk Blaschke
factors
\[
 B_a(z)=\frac{R(z-a)}{R^2-\overline a z},
\]
repeated with multiplicity.  The quotient $G$ is holomorphic and
nonvanishing on the closed disk, and has the same boundary modulus
as $D$.  Maximum modulus gives $\log|G|\le M_0$ throughout the
disk for some $M_0=O_C(T^2)$.  Since $|B_a(0)|<1$,
$\log|G(0)|\ge-O(T)$.

Let $h$ be an analytic logarithm of $G$.  The positive harmonic
function $M_0-\operatorname{Re}h$ has value $O_C(T^2)$ at the
center.  The Poisson formula, or its derivative together with Harnack's
inequality on concentric disks, gives
\[
 |h'(z)|\ll_C T^2/r_*\qquad(|z|\le2r_*).
\]
Restoring the factors uses
\[
 \frac{B_a'}{B_a}(z)=\frac1{z-a}
       +\frac{\overline a}{R^2-\overline a z}.
\]
The second term is $O(1/r_*)$ on $|z|\le2r_*$, uniformly in
$|a|<R$.  There are $O_C(T^2)$ such terms, giving exactly
\eqref{lp:local-log-derivative-formula}.
\end{proof}

\begin{lemma}[A zero-free strip at large height]\label{lp:zero-free}
For each fixed $C>0$ there is $c=c(C)>0$ such that every character
of modulus at most $L^C$ has no zero in
\begin{equation}\label{lp:zero-free-region}
 \operatorname{Re}s\ge1-cT^2/L,\qquad
                    1\le|\operatorname{Im}s|\le x^3.
\end{equation}
Moreover,
\begin{equation}\label{lp:log-derivative-bound}
 \left|\frac{D'}D(s,\chi)\right|\ll_C L
 \quad\left(1-\frac{cT^2}{2L}\le\operatorname{Re}s\le1+\frac1L,
       \quad2\le|\operatorname{Im}s|\le\frac{x^3}{2}\right).
\end{equation}
\end{lemma}

\begin{proof}
For $\sigma>1$, the Euler products and the inequality
$3+4\cos\theta+\cos(2\theta)=2(1+\cos\theta)^2\ge0$ give
\begin{equation}\label{lp:trigonometric-positivity}
 0\le-3\frac{\zeta'}\zeta(\sigma)
       -4\operatorname{Re}\frac{D'}D(\sigma+iv,\chi)
       -\operatorname{Re}\frac{D'}D(\sigma+2iv,\chi^2).
\end{equation}
Indeed this follows term by term in the absolutely convergent
prime-power expansions; primes dividing the modulus contribute only
the positive zeta term.  Also
$-\zeta'/\zeta(\sigma)=1/(\sigma-1)+O(1)$ near $1$.

Suppose $\rho=\beta+iv$ were a zero in
\eqref{lp:zero-free-region}, and set
$\sigma=1+20cT^2/L$.  Apply
Lemma~\ref{lp:local-log-derivative} at heights $v$ and $2v$.
The evaluation points lie in the respective inner disks, and $\rho$
lies in the first zero sum, because $T^2/L=o(r_*)$.
No zero has real part greater than $1$, by the absolutely convergent
Euler product.  All terms in the zero sums therefore contribute
nonpositively to the negative real logarithmic derivatives.  Retaining
the term at $\rho$, and using
$0<\sigma-\beta\le21cT^2/L$, bounds the right-hand side of
\eqref{lp:trigonometric-positivity} by
\[
 \left(\frac{3}{20c}-\frac{4}{21c}+O_C(1)\right)\frac L{T^2}
 =\left(-\frac{17}{420c}+O_C(1)\right)\frac L{T^2}.
\]
Here the errors are $O_C(T^2/r_*)=O_C(L/T^2)$.
Choosing a sufficiently small fixed $c>0$ gives a contradiction.

For $s$ in the region of \eqref{lp:log-derivative-bound}, apply the
local formula centered at $1+r_*+i\operatorname{Im}s$.
Every zero in its sum has absolute imaginary part between $1$ and
$x^3$, since $2\le|\operatorname{Im}s|\le x^3/2$ and $r_*=o(1)$.
By \eqref{lp:zero-free-region}, its horizontal distance from $s$ is
at least $cT^2/(2L)$.  The $O_C(T^2)$ zero terms and the local error
therefore total $O_C(L)$, proving
\eqref{lp:log-derivative-bound}.
\end{proof}

\begin{proof}[Proof of Lemma~\ref{lp:prime-polynomial}]
We first establish the analogous bound with the von Mangoldt weight.
Let $\delta=L^{-A-4}$.  Smooth the indicator of $I/N\subset[1,2]$
by convolution with a nonnegative smooth kernel of width $\delta$.
This gives $0\le g\le1$, supported in $[1/2,3]$, whose difference
from the indicator is supported within $O(\delta)$ of its endpoints,
and with $\|g^{(j)}\|_\infty\ll_j\delta^{-j}$.
The construction also applies when $I$ is shorter than $\delta N$.
Since $\Lambda(n)\le\log n$, the error in replacing the interval
by $g(n/N)$ is
\begin{equation}\label{lp:smoothing-error}
 O\!\left((\delta N+1)\frac{\log(3N)}N\right)
       =O(L^{-A-2}),
\end{equation}
uniformly in $I$.

Define the Mellin transform by
\[
 \widehat g(z)=\int_0^\infty g(w)w^z\,\frac{dw}{w}.
\]
On every fixed bounded real-part strip, integration by parts gives
\begin{equation}\label{lp:mellin-decay}
 |\widehat g(u+iv)|\ll_j
       L^{(j+1)(A+5)}(1+|v|)^{-j}.
\end{equation}
Also $|\widehat g(u+iv)|\ll1$ there, by absolute integration.
Mellin inversion and the absolutely convergent logarithmic derivative
on $\operatorname{Re}z=1/L$ yield
\begin{equation}\label{lp:mellin-inversion}
 \sum_n\Lambda(n)\chi(n)n^{-1+it}g(n/N)
 =\frac1{2\pi i}\int_{1/L-i\infty}^{1/L+i\infty}
       \widehat g(z)N^z
       \left(-\frac{D'}D(1-it+z,\chi)\right)dz.
\end{equation}
On this full line, absolute convergence gives
\[
 \left|\frac{D'}D(1+1/L+i u,\chi)\right|
 \le\sum_n\frac{\Lambda(n)}{n^{1+1/L}}
       =-\frac{\zeta'}\zeta(1+1/L)\ll L,
\]
at every height $u$.

Choose a fixed $C_5>A+6$ and put $H_0=L^{C_5}$.
Then choose a fixed integration-by-parts order $j$ large enough that
\eqref{lp:mellin-decay} makes the two tails with
$|\operatorname{Im}z|>H_0$ in \eqref{lp:mellin-inversion}
$O(L^{-A-2})$.  Explicitly their bound is
\[
 O_j\!\left(L^{1+(j+1)(A+5)}H_0^{1-j}\right),
\]
since $N^{1/L}\ll1$.  Fix $B_0>C_5+2$.
For $L^{B_0}\le|t|\le x^2$, every point in the rectangle
\[
 -\frac{cT^2}{2L}\le\operatorname{Re}z\le\frac1L,
       \qquad |\operatorname{Im}z|\le H_0
\]
has
\begin{equation}\label{lp:contour-heights}
 2\le|-t+\operatorname{Im}z|
       \le x^2+H_0<\frac{x^3}{2}
\end{equation}
for large $x$.  Thus Lemma~\ref{lp:zero-free} applies throughout
the rectangle.  There are no zeros or poles of the logarithmic
derivative inside it; in particular, the possible principal-character
pole at $z=it$ is outside it.

Shift the truncated contour to
$\operatorname{Re}z=-cT^2/(2L)$.
The horizontal edges are $O(L^{-A-2})$, by
\eqref{lp:log-derivative-bound} and \eqref{lp:mellin-decay}, increasing
the already fixed order $j$ if necessary.
On the new vertical segment,
\[
 |N^z|=\exp\left(-\frac{cT^2\log N}{2L}\right)
       \ll\exp(-c\tau T^2/2).
\]
Its remaining factors and length cost at most a fixed power of $L$.
Since $\exp(-c\tau T^2/2)$ is smaller than every prescribed fixed
power of $L^{-1}$, the shifted integral is $O(L^{-A-2})$.
Together with \eqref{lp:smoothing-error}, this proves, uniformly for
all subintervals $I\subset[N,2N]$,
\begin{equation}\label{lp:mangoldt-estimate}
 \sum_{n\in I}\Lambda(n)\chi(n)n^{-1+it}\ll L^{-A-2}.
\end{equation}

Prime powers of exponent at least two contribute in absolute value at
most $O(N^{-1/2}(\log(3N))^2)$: there are
$O(\sqrt N\log(3N))$ possible bases and exponents with
$N\le p^a\le2N$, and every summand has size $O(\log(3N)/N)$.
This is smaller than every fixed power of $L^{-1}$ in the present
range.  Removing them from \eqref{lp:mangoldt-estimate} gives the
same uniform bound for
$\sum_{p\in I}(\log p)\chi(p)p^{-1+it}$.
Finally, partial summation against $1/\log u$, whose value and total
variation on $[N,2N]$ are $O_{\tau}(1/L)$, removes the logarithmic
weight and proves \eqref{lp:prime-estimate}.
\end{proof}
